
@@M0@@
where @@M1@@, @@M2@@ are the \emph{geometric} fibers of @@M3@@. Moreover, one knows that the cohomological specialization homomorphism @@M4@@ is now an isomorphism (SGA 4 XVI 2.2), so the commutativity of the preceding square can also be expressed as the commutativity of the square obtained by reversing the preceding arrow. Passing to the projective limit over @@M5@@ in the commutative squares thus obtained, one finds a commutative square
\sgaStableTarget{sga6.e10.eq.7-13-4}@@M6@@

\sgaStableTarget{sga6.e10.para.7-13-5}\textbf{7.13.5.} Another way to express this commutativity is to say that, when @@M7@@ is an element of @@M8@@ and @@M9@@ is its image in @@M10@@, @@M11@@ and @@M12@@ are the values, at the geometric points @@M13@@ and @@M14@@ respectively, of a \emph{section} of @@M15@@. This section is moreover readily interpreted by taking an element @@M16@@ inducing @@M17@@, whence a @@M18@@, which defines the required section of @@M19@@.

\sgaStableTarget{sga6.e10.para.7-13-6}\textbf{7.13.6.} Let @@M20@@ be a scheme. Then, by passage to the associated graded groups, the homomorphism (7.13.0) defines homomorphisms
@@M21@@
the image under @@M22@@ of an element @@M23@@ of @@M24@@ being, by definition, @@M25@@. (The fact that (7.13.0) is compatible with the @@M26@@-filtration on
